\documentclass[journal,10pt,twocolumn]{IEEEtran}
\usepackage{amsmath,amssymb,amsfonts}
\usepackage{graphicx}
\usepackage{booktabs}
\usepackage{cite}
\usepackage{microtype}
\usepackage{hyperref}
\usepackage{xcolor}

\hypersetup{
    colorlinks=true,
    linkcolor=blue,
    citecolor=blue,
    urlcolor=blue
}

\begin{document}

\title{Context-Adaptive Expert Gating Network for Short-Term Electricity Load Forecasting with Causal Out-of-Fold Performance Feedback and Confidence Fallback}

\author{[Author Name(s)]\\
[Department / Institution]\\
[Corresponding Author Email]}

\maketitle

\begin{abstract}
Accurate short-term electricity load forecasting (STLF) over a 24-hour horizon is critical for power grid unit commitment, reserve margin scheduling, and market operations. While diverse deep sequence architectures (LSTM, TCN, CNN) offer complementary inductive biases, standard mixture-of-experts (MoE) gating networks suffer from overfitting, training-side target leakage, and unconstrained over-allocation under distribution shift. We present CAEG-Net, a context-adaptive expert gating framework integrating three canonical temporal experts: a 2-layer LSTM (56,152 parameters), a 6-stage dilated causal TCN with a 253-hour receptive field (36,952 parameters), and a 3-stage multi-scale CNN (27,400 parameters), totaling 120,504 frozen core parameters. CAEG-Net incorporates a 7-dimensional context representation coupling 4 operational state features (trend, volatility, lag-24 autocorrelation, recent error) with 3 chronological out-of-fold (OOF) relative expert-performance metrics generated via expanding-window cross-validation over 4 training blocks. A dedicated confidence head estimates an empirical balance parameter $\lambda_t \in (0, 1)$ blending dynamic routing with an uninformative equal-weight prior ($[1/3, 1/3, 1/3]$). Evaluated across Modern PJM (MW), GEFCom2014 (kW), and UCI Electricity (Cohort 320, MW) benchmarks over 5 evaluation seeds ($\{42, 123, 999, 2024, 3407\}$), CAEG-Net (Candidate \texttt{F2\_A2\_OOF}) achieves statistically significant reductions in MAE over canonical context-only gating across all three datasets under Holm-Bonferroni correction (Modern PJM: $\bar{\Delta} = -9.66$ MW, $p_{\mathrm{adj}} = 0.0406$; GEFCom: $\bar{\Delta} = -0.688$ kW, $p_{\mathrm{adj}} = 1.02 \times 10^{-27}$; UCI: $\bar{\Delta} = -0.202$ MW, $p_{\mathrm{adj}} = 0.0017$). Furthermore, CAEG-Net achieves lower MAE than static equal ensembling on all three benchmarks (PJM: $-10.31\%$, GEFCom: $-1.66\%$, UCI: $-5.26\%$) with only 193 additional parameters (+0.1588\%). The learned confidence coefficient exhibits low temporal variability ($\lambda \approx 0.51 \pm 0.005, CV < 1.1\%$) and behaves approximately as a near-constant shrinkage toward equal fusion rather than a rapid dynamic switching circuit.
\end{abstract}

\begin{IEEEkeywords}
Electricity load forecasting, Mixture of Experts, Adaptive gating, Out-of-fold performance feedback, Forecast combinations, Shrinkage methods.
\end{IEEEkeywords}

\section{Introduction}
Short-term electricity load forecasting (STLF) is an essential operational requirement for electrical power system security and day-ahead market efficiency. Grid dispatchers rely on 24-hour forward load curves to solve mixed-integer linear programming (MILP) unit commitment problems and maintain spinning reserve requirements. Modern power grids exhibit pronounced non-stationarity due to behind-the-meter photovoltaics, EV charging, and demand electrification.

While deep sequence models such as Long Short-Term Memory (LSTM)~\cite{hochreiter1997long}, Temporal Convolutional Networks (TCN)~\cite{bai2018empirical}, and multi-scale Convolutional Neural Networks (CNN) capture distinctive temporal dynamics, no single architecture uniformly dominates across all operational regimes. Static equal ensembling offers strong empirical variance reduction~\cite{clemen1989combining,timmermann2006forecast} but remains non-adaptive. Mixture-of-Experts (MoE)~\cite{jacobs1991adaptive} introduces dynamic gating, but often suffers from gating overfitting, in-sample training performance leakage, and routing instability.

This paper establishes CAEG-Net, combining heterogeneous temporal experts with causally constructed out-of-fold performance feedback and a confidence-guided shrinkage fallback, validated under strict chronological evaluation across three heterogeneous benchmarks.

\section{Methodology and Architecture}
\subsection{Problem Formulation}
Given historical sequence $X_{t-L+1:t} \in \mathbb{R}^L$ ($L=168$), the model predicts $Y_{t+1:t+H} \in \mathbb{R}^H$ ($H=24$). Optimization minimizes Mean Squared Error (MSE), while Mean Absolute Error (MAE) serves as the primary evaluation metric:
\begin{equation}
\text{MAE} = \frac{1}{NH} \sum_{i=1}^N \sum_{h=1}^H |y_{t_i+h} - \hat{y}_{t_i+h}|
\end{equation}

\subsection{CAEG-Net Architecture}
CAEG-Net incorporates three canonical temporal experts:
\begin{itemize}
    \item \textbf{LSTM Expert:} 2-layer LSTM ($h=64$, dropout $0.1$) with forecasting head Linear(64, 64) $\to$ ReLU $\to$ Linear(64, 24) (56,152 params).
    \item \textbf{TCN Expert:} 6 causal residual stages ($k=3$, dilations $\{1, 2, 4, 8, 16, 32\}$, 32 ch, RF=253h) with forecasting head Linear(32, 32) $\to$ ReLU $\to$ Linear(32, 24) (36,952 params).
    \item \textbf{CNN Expert:} 3-stage Conv1D ($1 \to 32 \to 64 \to 64$ with MaxPool and AdaptiveAvgPool) with forecasting head Linear(64, 48) $\to$ ReLU $\to$ Linear(48, 24) (27,400 params).
\end{itemize}
The expert core comprises 120,504 parameters.

\subsection{Chronological OOF Performance Feedback}
To avoid training-side target leakage, expanding-window cross-validation is performed across 4 chronological training blocks. For expert $i$, historical error is $e_i(t) = \frac{1}{24} \sum_{h=1}^{24} |y_{t-h} - \hat{y}_{i, t-h}|$, yielding relative error $r_i(t) = e_i(t) / \sum_j e_j(t)$. The 3 relative error signals are concatenated with 4 operational state features (trend, volatility, lag-24 autocorrelation, recent error) into a 7-dimensional context vector $c_t \in \mathbb{R}^7$.

\subsection{Confidence Fallback Head}
A softmax router predicts expert weights $w_t \in \Delta^2$ (1,075 params), while a confidence head predicts balance parameter $\lambda_t \in (0, 1)$ (145 params):
\begin{equation}
\hat{y}_{t}^{\mathrm{final}} = \lambda_t \sum_{m=1}^3 w_{t, m} \hat{y}_{t}^{(m)} + (1 - \lambda_t) \frac{1}{3} \sum_{m=1}^3 \hat{y}_{t}^{(m)}
\end{equation}
The gating and confidence heads add 1,220 parameters (net overhead +193 parameters, +0.1588\% over Canonical V1).

\section{Experimental Results}
Table~\ref{tab:results} reports five-seed held-out test performance (Mean $\pm$ Population SD, \texttt{ddof=0}) across benchmarks.

\begin{table}[htbp]
\caption{Held-Out Test MAE Across Five Random Seeds}
\label{tab:results}
\centering
\begin{tabular}{lccc}
\toprule
Model & Modern PJM & GEFCom2014 & UCI Cohort 320 \\
 & (MW) & (kW) & (MW) \\
\midrule
Static Equal Ens. & 279.83 & 12.62 & 8.17 \\
Best Standalone & 259.33 (TCN) & 12.57 (TCN) & 7.55 (LSTM) \\
Canonical V1 & $253.41 \pm 9.12$ & $12.88 \pm 0.25$ & $7.94 \pm 0.17$ \\
\textbf{CAEG-Net (F2)} & $\mathbf{250.97 \pm 10.69}$ & $\mathbf{12.41 \pm 0.15}$ & $\mathbf{7.74 \pm 0.30}$ \\
\bottomrule
\end{tabular}
\end{table}

\section{Statistical Inference}
Hypothesis testing on non-overlapping 24-hour daily blocks ($K=53, 456, 163$) shows significant reductions in MAE over Canonical V1 under Holm-Bonferroni correction:
\begin{itemize}
    \item \textbf{PJM ($K=53$):} $\bar{\Delta} = -9.66$ MW, $p_{\mathrm{adj}} = 0.0406$, Cohen's $d_z = -0.351$.
    \item \textbf{GEFCom ($K=456$):} $\bar{\Delta} = -0.688$ kW, $p_{\mathrm{adj}} = 1.02 \times 10^{-27}$, Cohen's $d_z = -0.555$.
    \item \textbf{UCI ($K=163$):} $\bar{\Delta} = -0.202$ MW, $p_{\mathrm{adj}} = 0.0017$, Cohen's $d_z = -0.287$.
\end{itemize}

\section{Conclusion}
This paper presented CAEG-Net, a context-adaptive expert gating framework for short-term electricity load forecasting. The final \texttt{F2\_A2\_OOF} formulation combines heterogeneous LSTM, TCN, and CNN experts with chronologically constructed out-of-fold expert-performance features and a confidence-guided fallback toward equal fusion. Across Modern PJM, GEFCom2014, and the UCI Cohort 320 benchmark, F2 improved upon canonical V1 in the locked five-seed evaluation and achieved lower MAE than static equal fusion on all three datasets. Under the predefined non-overlapping daily-block paired analysis with Holm correction, F2 showed statistically significant improvement over V1 across all three datasets. F2 was not the lowest-error formulation on every individual dataset and did not outperform the strongest standalone expert on UCI under the locked reference. Its selection therefore reflects cross-dataset robustness and consistency rather than universal per-dataset optimality. The confidence coefficient exhibited low temporal variability, suggesting that the fallback mechanism behaved approximately as a conservative shrinkage toward equal fusion rather than as a strongly time-varying expert switch. Overall, the results provide empirical evidence that causally constructed performance-aware adaptive fusion can improve the robustness of heterogeneous deep forecasting ensembles under strict chronological evaluation.

\bibliographystyle{IEEEtran}
\bibliography{references/references}

\end{document}
